Сближение России и Великобритании после 1907 года показывает, что отношения между государствами не являются неизменными: многолетняя вражда может смениться сотрудничеством при изменении внешних условий, например усиление общего врага. Это подтверждает, что в мировой политике нет «вечных союзников» или «вечных врагов», а есть прежде всего интересы. Поэтому международные отношения нельзя рассматривать в категориях «чёрного» и «белого». Она гораздо сложнее и ближе к оттенкам серого, где каждое решение — это результат компромиссов, расчёта и изменяющихся условий. Необходимо понимать это при анализе исторических процессов или политической ситуации в мире.

На этом - всё.
Рекомендую взять подмышку любимых и сходить куда-нибудь вместе. Этот момент вы точно запомните в отличии от этого ролика. Всем хорошего дня.